\section{Medical and Healthcare Explainability Strategy}

Because \projectname{} is intended for medical and healthcare data workflows, explainability must be treated as a core system requirement. The system should not merely ask a foundation model to ``explain itself.'' In clinical settings, natural-language rationales can be persuasive but unfaithful; recent work shows LLM self-explanations should not be trusted by default \cite{huang2024selfExplanations}. Explanation must therefore be built from verifiable evidence, validated data operations, uncertainty estimates, human review, and audit trails.

\subsection{Explainability Position}

\projectname{} should use a layered explainability model:

\begin{itemize}
  \item \rolelabel{Data explainability:} what data was used, what format it had, what schema was inferred, what fields were missing, what values looked anomalous, and what transformations were applied.
  \item \rolelabel{Retrieval explainability:} which clinical guideline, schema, ontology entry, data dictionary, or prior validated example was retrieved; why it was considered relevant; and what evidence was not found.
  \item \rolelabel{Document and visual explainability:} which OCR page, layout block, table cell, DICOM study, image slice, frame, mask, box, or tracking segment contributed to a result.
  \item \rolelabel{Model explainability:} what prompt, retrieved evidence, tool outputs, and uncertainty signals shaped the final response. The system should not present chain-of-thought as a faithful internal explanation.
  \item \rolelabel{Workflow explainability:} which agents acted, which MCP tools were called, what validation checks passed or failed, and when human approval was required.
  \item \rolelabel{Clinical explainability:} a clinician-facing summary of intended use, limitations, uncertainty, data-quality warnings, and recommended human review points.
\end{itemize}

\begin{figure}[H]
  \centering
  \resizebox{\textwidth}{!}{\begin{tikzpicture}[
  font=\small,
  >=Latex,
  box/.style={draw, rounded corners=4pt, minimum height=0.78cm, align=center, inner sep=5pt},
  layer/.style={box, text width=2.6cm},
  arrow/.style={-{Latex[length=2mm]}, thick}
]

\node[layer, fill=sfblue!8, draw=sfblue] (input) {Healthcare data\\EHR, labs, CSV, FHIR-like records};
\node[layer, fill=sfcyan!8, draw=sfcyan, right=0.55cm of input] (validate) {Validation trace\\schema, units, missingness};
\node[layer, fill=sfgreen!8, draw=sfgreen, right=0.55cm of validate] (retrieval) {Evidence trace\\guidelines, schemas, dictionaries};
\node[layer, fill=sfpurple!8, draw=sfpurple, right=0.55cm of retrieval] (uncertainty) {Uncertainty\\confidence, conflicts, abstention};
\node[layer, fill=sforange!10, draw=sforange, right=0.55cm of uncertainty] (clinician) {Clinical review\\approve, edit, reject};
\node[layer, fill=sfblue!8, draw=sfblue, right=0.55cm of clinician] (explain) {Explanation report\\claims, sources, limits};

\draw[arrow, sfblue] (input) -- (validate);
\draw[arrow, sfcyan] (validate) -- (retrieval);
\draw[arrow, sfgreen] (retrieval) -- (uncertainty);
\draw[arrow, sfpurple] (uncertainty) -- (clinician);
\draw[arrow, sforange] (clinician) -- (explain);

\node[box, fill=sfgray, draw=sfdark, text width=3.1cm, below=0.75cm of retrieval] (audit) {Audit log\\workflow ID, tool calls, source IDs, approvals};
\node[box, fill=sfred!7, draw=sfred, text width=3.2cm, below=0.75cm of uncertainty] (guard) {Guardrail\\not diagnosis, not treatment, escalate when unsure};

\draw[arrow, sfdark] (validate) |- (audit);
\draw[arrow, sfdark] (retrieval) -- (audit);
\draw[arrow, sfdark] (clinician) |- (audit);
\draw[arrow, sfred] (guard) -- (uncertainty);

\end{tikzpicture}
}
  \caption{Healthcare-grade explainability stack for OJTFlow. Explanations are assembled from evidence, validation, uncertainty, and review traces rather than unsupported model rationales.}
  \label{fig:medical-explainability}
\end{figure}

\subsection{SOTA Explainability Methods to Apply}

\begin{longtable}{p{0.19\textwidth}p{0.32\textwidth}p{0.32\textwidth}}
\toprule
\textbf{Method} & \textbf{Purpose} & \textbf{Use in OJTFlow} \\
\midrule
Evidence-grounded RAG explanations & Retrieval provides inspectable source evidence and provenance, reducing reliance on parametric model memory. & Every clinical or healthcare explanation should include source IDs, schema versions, retrieved context summaries, and missing-evidence warnings. \\
\addlinespace
Citation checks & Generated claims should be checked against retrieved evidence and validation reports. & The Explanation Agent should mark claims as \texttt{supported}, \texttt{partially supported}, \texttt{unsupported}, or \texttt{requires clinician review}. \\
\addlinespace
Counterfactual and contrastive explanations & Explain why one schema, field mapping, or transformation was chosen instead of another. & Show alternative schema candidates and the evidence that favored the selected schema. \\
\addlinespace
Feature attribution & LIME, SHAP, Integrated Gradients, and TCAV provide local or concept-level explanations for model components \cite{ribeiro2016lime,lundberg2017shap,sundararajan2017ig,kim2018tcav}. & Use for auxiliary classifiers such as PII detection, anomaly classification, schema matching, or risk scoring. Do not overclaim them as full LLM explanations. \\
\addlinespace
Attention and saliency with caveats & Attention weights are often tempting as explanations, but NLP research warns that attention is not necessarily a faithful explanation \cite{jain2019attention}. & If attention/saliency is shown, label it as a diagnostic visualization, not proof of clinical reasoning. \\
\addlinespace
Mechanistic and probing methods & LLM explainability surveys categorize approaches such as probing, attribution, neuron/activation analysis, and mechanistic interpretation \cite{zhao2024llmExplainability}. & Future research track for model behavior auditing, not a required MVP dependency. \\
\addlinespace
Uncertainty, calibration, and abstention & Medical LLM evaluation emphasizes reliability, calibration, and knowing when not to answer \cite{chen2024medEval}. & Require confidence scores, retrieval sufficiency checks, and explicit abstention/clinician escalation when evidence is incomplete. \\
\addlinespace
Human-centered XAI evaluation & Clinical AI reporting and evaluation should include safety, human factors, workflow fit, and stakeholder-specific explanations \cite{vasey2022decide,fdaTransparency}. & Evaluate explanations with reviewers using clarity, actionability, trust calibration, and risk-awareness rubrics. \\
\addlinespace
Visual evidence explanation & Medical OCR, segmentation, detection, and tracking should be explained with source page regions, image overlays, uncertainty, and reviewer edits rather than text-only rationales. & Show OCR bounding boxes, masks, boxes, frame IDs, image metadata, and clinician correction history for every multimodal output. \\
\bottomrule
\end{longtable}

\subsection{Clinical Explanation Output Contract}

For medical and healthcare workflows, final responses should include a structured explanation object:

\begin{verbatim}
{
  "answer_type": "data_transformation_explanation",
  "intended_use": "Support data validation and review; not autonomous diagnosis",
  "evidence": [
    {"source_id": "schema:lab_result_v2", "claim": "LOINC code is required"},
    {"source_id": "validation:wf_001", "claim": "3 rows have missing units"}
  ],
  "data_quality_flags": ["missing_units", "date_format_inconsistency"],
  "uncertainty": {
    "retrieval_confidence": "medium",
    "schema_confidence": 0.78,
    "requires_clinician_review": true
  },
  "transformation_trace": ["parse_csv", "infer_schema", "validate_schema"],
  "multimodal_evidence": {
    "ocr_boxes": ["page_2:block_7"],
    "image_regions": ["study_01:series_02:slice_18:mask_3"],
    "video_tracks": []
  },
  "unsupported_claims": [],
  "limitations": ["No diagnosis or treatment recommendation was generated"]
}
\end{verbatim}

\subsection{Healthcare Safety and Governance Rules}

\begin{itemize}
  \item The system should not make autonomous diagnosis, treatment, triage, or medication recommendations in the MVP.
  \item Explanations must separate data facts, retrieved source statements, model-generated summaries, and human decisions.
  \item Every clinically relevant claim should be traceable to source evidence, a validation result, or a human approval.
  \item Every multimodal claim should be traceable to a page/box, table cell, DICOM study/series/instance, image slice, mask, bounding box, frame, or reviewer correction.
  \item Low retrieval confidence, conflicting sources, missing patient context, or schema ambiguity should trigger abstention or clinician review.
  \item Patient-facing explanations should use plain language; clinician-facing explanations may include technical terminology, coding systems, confidence intervals, and validation details.
  \item Explanations should include known limitations, data gaps, and possible bias or distribution-shift warnings.
\end{itemize}

\subsection{Medical Evaluation Plan}

\begin{tabularx}{\textwidth}{p{0.26\textwidth}X p{0.23\textwidth}}
\toprule
\textbf{Evaluation area} & \textbf{Metric or test} & \textbf{Target} \\
\midrule
Evidence faithfulness & Percentage of explanation claims supported by retrieved sources, validation reports, or tool outputs. & \metric{>= 95\% supported for demo cases}. \\
Clinical clarity & Reviewer rating for whether explanation is understandable and useful for the intended audience. & \metric{Average >= 4/5}. \\
Uncertainty handling & Correct abstention/escalation on incomplete, conflicting, or low-confidence evidence. & \metric{No unsafe confident answer}. \\
Human factors & Reviewer can identify what changed, why it changed, and whether approval is required. & \metric{>= 90\% task success}. \\
Bias and data gaps & Presence of warnings for missing demographics, skewed samples, unknown units, and schema mismatch. & \metric{All seeded gaps flagged}. \\
Auditability & Ability to reconstruct input, retrieval, tool calls, validation, approval, and final answer. & \metric{100\% for completed workflows}. \\
\bottomrule
\end{tabularx}
